% Davis-Kahan sin-Theta: what the Lean formalization proves, and how it relates
% to the literature.  Written as an expert-readable companion to
% ForMathlib/Analysis/InnerProductSpace/DavisKahan.lean, so that the statement
% can be checked against the source papers without reading Lean.
%
% Independently written; the two source-paper digests it draws on are
%   ForMathlib/prose/Davis-1963-core-arguments.tex
%   ForMathlib/prose/Yu-Wang-Samworth-2014-core-arguments.tex

\documentclass[11pt]{article}

\usepackage[T1]{fontenc}
\usepackage{lmodern}
\usepackage{microtype}
\usepackage[margin=1.1in]{geometry}
\usepackage{amsmath,amssymb,amsthm,mathtools}
\usepackage{booktabs}
\usepackage[colorlinks=true, urlcolor=blue, citecolor=blue, linkcolor=black]{hyperref}

\newcommand{\norm}[1]{\left\lVert #1\right\rVert}
\newcommand{\eps}{\varepsilon}
\newcommand{\F}{\mathrm{F}}
\newcommand{\op}{\mathrm{op}}
\newcommand{\inner}[2]{\langle #1,\,#2\rangle}
\newcommand{\gap}{g}
\newcommand{\HH}{\mathcal H}
\newcommand{\sinT}{\sin\Theta}
\newcommand{\tp}{^{\mathsf T}}
\DeclareMathOperator{\tr}{tr}
\DeclareMathOperator{\ran}{ran}
\DeclareMathOperator{\spann}{span}

\newtheorem{theorem}{Theorem}
\newtheorem{corollary}[theorem]{Corollary}
\newtheorem{lemma}[theorem]{Lemma}
\theoremstyle{remark}
\newtheorem*{remark}{Remark}

\title{Davis--Kahan $\sin\Theta$: what the Lean formalization proves,\\
and how it relates to Davis (1963/1970) and Yu--Wang--Samworth (2015)}

% NOTE: if you are a different model making a substantial addition and are not
% yet credited, add yourself to the author list above.
\author{Claude Opus 4.8 \and Jonathan Crall}

\date{\today}

\begin{document}
\maketitle

\begin{abstract}
This note states, in ordinary mathematical language, the Davis--Kahan (DK)
$\sin\Theta$ theorem formalized in Lean, and locates it among the source papers.
The formalization is a ladder of Frobenius-norm $\sin\Theta$ bounds built on one
estimate, whose sharpest rung is $\norm{\sinT}_{\F}\le\norm{S-T}_{\F}/\gap$ with a
dimension-free constant.  The estimate is Davis's elementary cross-term /
Rayleigh-quotient inequality, and the separation is a hybrid gap between the
leading eigenvalues of one operator and the trailing eigenvalues of the other.
The hybrid-gap bound is the deterministic core of the Yu--Wang--Samworth (YWS)
variant; recovering their population-only gap is a further Weyl step, also
formalized.  A closing section records the DK variants left unformalized and what
each of them requires.
\end{abstract}

\vspace{1.5em}
% NOTE: Update permalink if underlying proof changes, so this always is up to date wrt to the referenced lean file.
\noindent Companion note to \href{https://github.com/AIQ-Kitware/aiq-dkps-formalization/blob/375251df7ab712397977635e8f607c11c3a8083d/ForMathlib/Analysis/InnerProductSpace/DavisKahan.lean}{ForMathlib/Analysis/InnerProductSpace/DavisKahan.lean}, part of a larger formalization effort at \href{https://github.com/AIQ-Kitware/aiq-dkps-formalization}{github.com/AIQ-Kitware/aiq-dkps-formalization}.



\section{Prerequisites}

This note assumes basic familiarity with the following finite-dimensional linear
algebra concepts.

\begin{itemize}
\item \textbf{Inner product spaces.}
We work over $\mathbb R$ or $\mathbb C$.  Inner products let us define
orthogonality, norms, orthonormal bases, and orthogonal projections.

\item \textbf{Self-adjoint operators.}
An operator $T$ is self-adjoint if $T=T^*$.  For complex matrices in an
orthonormal basis, this means $T$ equals its conjugate transpose; over
$\mathbb R$, this is the usual notion of a symmetric matrix.  Self-adjoint
operators have real eigenvalues and orthonormal eigenbases.

\item \textbf{Eigenvectors, eigenspaces, and spectral subspaces.}
An eigenspace is the span of all eigenvectors associated with a given
eigenvalue.  A spectral subspace is the span of eigenvectors associated with a
chosen group of eigenvalues, such as the largest $d$ eigenvalues.

\item \textbf{Orthogonal projections.}
Every subspace $U$ of a finite-dimensional inner product space has an
orthogonal projection $P_U$.  Comparing subspaces can often be reduced to
comparing their projection operators.

\item \textbf{Spectral gaps.}
A spectral gap is a positive separation between a group of eigenvalues of
interest and the rest of the spectrum.  For ordered eigenvalues
$\lambda_1\geq\lambda_2\geq\cdots\geq\lambda_n$, the leading-$d$ gap is often
$\lambda_d-\lambda_{d+1}$.

\item \textbf{Perturbation size.}
If $S$ is a perturbed version of $T$, the perturbation is $E=S-T$.  Its size is
measured using a norm, such as the operator norm or the
Frobenius/Hilbert--Schmidt norm.

\item \textbf{Frobenius/Hilbert--Schmidt norm.}
For matrices, the Frobenius norm is the square root of the sum of squared
entry magnitudes.  For operators on finite-dimensional inner product spaces,
the analogous basis-independent norm is the Hilbert--Schmidt norm.
\end{itemize}


\section{Introduction to Davis--Kahan theory}

Davis--Kahan theory is about the stability of eigenspaces under perturbation.
Suppose $T$ is a self-adjoint operator and $S=T+E$ is a nearby self-adjoint
operator. The eigenvalues and eigenvectors of $S$ may differ from those of $T$,
but if the perturbation $E=S-T$ is small, we would like to know whether the
important eigenspaces of $S$ remain close to the corresponding eigenspaces of
$T$.

This question appears naturally in principal component analysis. For example,
$T$ might be an ideal or population covariance operator, while $S$ is an
empirical covariance operator estimated from finite data. The leading
eigenvectors of $T$ define the true principal subspace, while the leading
eigenvectors of $S$ define the empirical principal subspace. Davis--Kahan theory
gives deterministic bounds on how far the empirical subspace can be from the true
subspace.

The main qualitative message is simple: eigenspaces are stable when the
perturbation is small compared to the relevant spectral gap. More explicitly,
the amount the eigenspace moves is bounded, up to constants, by the size of the
perturbation divided by the spectral gap. Thus a small perturbation causes only
a small rotation when the eigenvalues defining the subspace are well separated
from the rest of the spectrum. If the gap is tiny, however, even a small
perturbation can cause a large change, because nearby eigenspaces can mix or
exchange roles.

It is important that Davis--Kahan theory compares subspaces, not individual
eigenvectors. Individual eigenvectors are not always stable or canonical. If an
eigenvalue has multiplicity greater than one, then its eigenspace has many
possible orthonormal bases, and no single basis is preferred. Even for a
one-dimensional eigenspace, an eigenvector is only determined up to sign over
$\mathbb R$, or up to phase over $\mathbb C$. The eigenspace itself is the
stable geometric object.

There are several equivalent ways to measure the distance between two subspaces.
One way is through principal angles. If $U$ and $V$ are two $d$-dimensional
subspaces, their principal angles $\theta_1,\dots,\theta_d$ measure how far the
directions in one subspace are from the directions in the other. The notation
$\sin\Theta$ refers to the collection of numbers
$\sin\theta_1,\dots,\sin\theta_d$. A norm such as
$|\sin\Theta|_{\mathrm F}$ then measures the total misalignment between the two
subspaces.

Another way to measure subspace distance is through orthogonal projections. Let
$P_U$ and $P_V$ be the orthogonal projection operators onto $U$ and $V$. If the
subspaces are close, then the projections $P_U$ and $P_V$ are close as
operators. Many versions of Davis--Kahan theory can therefore be stated either
in terms of principal angles or in terms of projection differences such as
$|P_U-P_V|_{\mathrm F}$.

A third way is through basis overlaps. Suppose $U$ is spanned by orthonormal
vectors $u_1,\dots,u_d$, and $V$ is spanned by orthonormal vectors
$v_1,\dots,v_d$. To measure how far $U$ is from $V$, one can measure how much of
the $u_i$ directions lie outside $V$. If
$v_{d+1},\dots,v_n$ span the orthogonal complement of $V$, then the quantity
\begin{equation}
\sum_{i=1}^d \sum_{j=d+1}^n |\langle u_i, v_j\rangle|^2
\end{equation}
measures how much the $U$-subspace leaks into the complement of the
$V$-subspace. This is another expression of the same $\sin\Theta$ geometry.
Here ``trailing'' means the complementary block $v_{d+1},\dots,v_n$, not another
block of size $d$.

The current Lean formalization is closest to this basis-overlap viewpoint, with
additional lemmas connecting it to projection distance. It works in a
finite-dimensional inner product space over $\mathbb R$ or $\mathbb C$, for
self-adjoint operators $T$ and $S$. The formalized theorem fixes a block size
$d$ and compares the span of the first $d$ eigenvectors of one operator with the
span of the first $d$ eigenvectors of the other. The theorem itself does not
require the eigenvalues to be globally sorted; it only requires an explicit gap
hypothesis separating the chosen block from its complement. In the usual
application, however, the chosen block is the leading eigenspace associated with
the largest $d$ eigenvalues.

The main formalized estimate is a Frobenius/Hilbert--Schmidt version of the
Davis--Kahan $\sin\Theta$ bound. It proves that a spectral gap controls the
cross-overlap between the chosen eigenspace of one operator and the complementary
eigenspace of the other. The proof then bounds this cross-overlap by
successively larger and easier-to-use quantities: first by an off-diagonal block
of the perturbation, then by a one-sided residual, and finally by the full
Hilbert--Schmidt norm of $S-T$. Additional lemmas use Weyl-type eigenvalue
perturbation bounds to turn more standard population-gap assumptions into the
hybrid gap assumption used by the main theorem.

This means the current formalization contains a substantial Frobenius-norm core
of Davis--Kahan theory, but it is not yet a complete formalization of all the
forms found in the literature. In particular, it does not yet provide a full
canonical-angle API, the dimension-free operator-norm $\sin\Theta$ theorem,
general unitarily invariant norm versions, arbitrary spectral intervals, the
aligned-basis conclusion from Yu--Wang--Samworth, or the full rotation theory
from Davis's 1963 paper. The purpose of this note is to explain exactly which
part has been formalized, how it relates to the paper statements, and what still
remains to be built.




\section{Setup and notation}

Let $\mathcal H$ be a finite-dimensional inner product space over
$\mathbb K\in{\mathbb R,\mathbb C}$ with $\dim\mathcal H=n$, and let
$T,S:\mathcal H\to\mathcal H$ be self-adjoint.  Write the eigenvalues and an
orthonormal eigenbasis of $T$ as $\lambda_1(T),\dots,\lambda_n(T)$ and
$u_1,\dots,u_n$, and those of $S$ as $\lambda_1(S),\dots,\lambda_n(S)$ and
$v_1,\dots,v_n$.  Fix a block size $d$.  In the intended top-eigenspace
application, the indices $1,\dots,d$ denote the leading eigenvectors, while
$d+1,\dots,n$ denote the complementary trailing eigenspace.  The formalized
hybrid-gap statement itself, however, only uses the block split
${1,\dots,d}$ versus ${d+1,\dots,n}$ together with an explicit separation
hypothesis between the two blocks; it does not require the eigenvalues to be
globally sorted unless one wants to derive that separation hypothesis from a
standard leading-eigenvalue gap.


Let $P$, $\widehat P$, and $Q=I-\widehat P$ be the orthogonal projections onto,
respectively,
\[
  \spann(u_1,\dots,u_d),\qquad
  \spann(v_1,\dots,v_d),\qquad
  \spann(v_{d+1},\dots,v_n).
\]
The first two are the leading spectral subspaces of $T$ and $S$, and $\ran Q$ is
the trailing $S$-subspace.  Let $\Theta$ be the vector of principal angles between
$\ran P$ and $\ran\widehat P$.  Two norms appear: the operator norm
$\norm{\cdot}_{\op}$ and the Frobenius, or Hilbert--Schmidt, norm
$\norm{B}_{\F}^2=\tr(B^*B)=\sum_k\norm{Be_k}^2$ over any orthonormal basis $e_k$.

The separation hypothesis is a \emph{hybrid gap}: a number $\gap>0$ with
\begin{equation}\label{eq:gap}
  \gap\le\bigl|\lambda_i(T)-\lambda_j(S)\bigr|
  \qquad\text{for all }i\le d<j.
\end{equation}
It pairs the leading eigenvalues of $T$ with the trailing eigenvalues of $S$, the
separation that Davis's estimate consumes.

\section{The formalized theorem}

The overlap $\sum_{i\le d}\sum_{j>d}|\inner{u_i}{v_j}|^2$ equals
$\norm{\sinT}_{\F}^2$, and the bounds on it form one chain of inequalities.

\begin{theorem}[Frobenius $\sin\Theta$ ladder]\label{thm:ladder}
Let $T,S$ be self-adjoint and let $\gap>0$ satisfy \eqref{eq:gap}.  Then
\begin{equation}\label{eq:ladder}
  \gap^2\,\norm{\sinT}_{\F}^2
  \;\le\;\norm{P(S-T)Q}_{\F}^2
  \;\le\;\norm{(S-T)P}_{\F}^2
  \;\le\;\norm{S-T}_{\F}^2,
\end{equation}
and consequently
\[
  \norm{\sinT}_{\F}\le\frac{\norm{S-T}_{\F}}{\gap},
  \qquad
  \norm{\widehat P-P}_{\F}\le\frac{\sqrt2\,\norm{S-T}_{\F}}{\gap},
\]
with constants independent of $n$.
\end{theorem}

The three right-hand sides of \eqref{eq:ladder} are the leading$\times$trailing
block of the perturbation, its restriction to the leading subspace (the one-sided
residual of the original DK theorem), and its full Frobenius norm; dividing by
$\gap^2$ gives three overlap bounds, of which the last is sharp.  The projector
distance and the overlap differ by an exact factor,
$\norm{\widehat P-P}_{\F}^2=2\norm{\sinT}_{\F}^2$, because
$\widehat P-P=\widehat P(I-P)-(I-\widehat P)P$ splits into two pieces with
orthogonal ranges of equal Frobenius norm, each equal to the overlap.

\begin{corollary}[operator-norm form]\label{cor:opnorm}
If $\norm{S-T}_{\op}\le\eps$, the residual has $d$ nonzero columns of norm at most
$\eps$, so $\norm{(S-T)P}_{\F}^2\le d\,\eps^2$ and
$\norm{\sinT}_{\F}^2\le d\,\eps^2/\gap^2$.  Extending to all $n$ columns gives the
cruder $\norm{\sinT}_{\F}^2\le n\,\eps^2/\gap^2$ and
$\norm{\widehat P-P}_{\F}^2\le 2n\,\eps^2/\gap^2$.  The example $S-T=\eps I$ shows
the dimension factor is unavoidable once only $\norm{S-T}_{\op}$ is controlled.
\end{corollary}

\begin{corollary}[rank and eigenvalue floor: the statistical setup]\label{cor:floor}
If $T$ is positive semidefinite of rank $d$ with nonzero eigenvalues at least
$\alpha$, and $\norm{S-T}_{\op}\le\eps\le\alpha/2$, then Weyl's inequality places
every trailing eigenvalue of $S$ below $\alpha/2$, so \eqref{eq:gap} holds with
$\gap=\alpha/2$ and $\norm{\sinT}_{\F}^2\le 4n\,\eps^2/\alpha^2$.
\end{corollary}

\section{The proof mechanism}

For eigenvectors $u_i$ of $T$ and $v_j$ of $S$, self-adjointness gives the
\emph{cross-term identity}
\begin{equation}\label{eq:cross}
  \inner{u_i}{(S-T)v_j}=\bigl(\lambda_j(S)-\lambda_i(T)\bigr)\inner{u_i}{v_j}.
\end{equation}
On a cross pair ($i\le d<j$), \eqref{eq:gap} turns this into
$\gap^2\,|\inner{u_i}{v_j}|^2\le|\inner{u_i}{(S-T)v_j}|^2$.  Summing over the
cross block gives the first inequality of \eqref{eq:ladder}, with right-hand side
$\norm{P(S-T)Q}_{\F}^2$; enlarging the column sum and then the row sum, by
Parseval in the bases $\{v_j\}$ and $\{u_i\}$, produces the remaining two rungs
and the full $\norm{S-T}_{\F}^2$.  The argument is elementary throughout, with no
resolvents or contour integrals.

\section{Relation to Davis (1963/1970)}

The estimate is Davis's~\cite{Davis1963}.  For a simple spectrum his
total-rotation bound reads, with $A\rightsquigarrow T$, $A+H\rightsquigarrow S$,
and $\mathcal C^{\perp}U$ the off-diagonal part of the matching unitary,
\[
  (\gamma')^2\,\norm{\mathcal C^{\perp}U}_{\F}^2
  \le\norm{H}_{\F}^2-\sum_i\bigl(\lambda_i-\lambda_i'\bigr)^2,
  \qquad
  \gamma_i=\min_{j\ne i}|\lambda_i-\lambda_j'|,
\]
proved by comparing $\inner{(A+H-\lambda_i')^2x_i}{x_i}$ from two sides, which is
\eqref{eq:cross}.  The formalization runs this argument on the two-block split.
The ladder \eqref{eq:ladder} realizes Davis's off-diagonal principle: the
cross-block rung $\norm{P(S-T)Q}_{\F}^2$ and the residual rung
$\norm{(S-T)P}_{\F}^2$ control the rotation by the block-mixing part of $H$,
and the sharp rung relaxes this to $\norm{S-T}_{\F}^2$.  The hybrid gap
\eqref{eq:gap} is Davis's $\gamma_i$ restricted to the cross block.  The companion
part~III of Davis and Kahan~\cite{DavisKahan1970} states the $\sin\Theta$ theorem
for every unitarily invariant norm; the Frobenius case is the instance formalized
here.

\section{Relation to Yu--Wang--Samworth (2015)}

Yu, Wang and Samworth~\cite{YWS2015} work with symmetric $\Sigma,\widehat\Sigma$
and a population eigengap $\Delta=\min(\lambda_{r-1}-\lambda_r,\
\lambda_s-\lambda_{s+1})$ formed from $\Sigma$ alone, and prove
\[
  \norm{\sinT(\widehat V,V)}_{\F}
  \le\frac{2\min\{d^{1/2}\norm{\widehat\Sigma-\Sigma}_{\op},
                   \norm{\widehat\Sigma-\Sigma}_{\F}\}}{\Delta}.
\]
Their proof is a residual sandwich
$\Delta\norm{\sinT}_{\F}\le\norm{R}_{\F}\le
2\min\{d^{1/2}\norm{\cdot}_{\op},\norm{\cdot}_{\F}\}$ with
$R=\widehat V\Lambda-\Sigma\widehat V$: a deterministic separation estimate on the
left, a perturbation estimate (Weyl and Hoffman--Wielandt) on the right, and the
factor $2$ paid for moving to the population-only gap $\Delta$.

The formalization keeps these two pieces separate.  The hybrid-gap bound, the
sharp rung of \eqref{eq:ladder}, is the deterministic separation estimate with
constant $1$: the hybrid gap already places the trailing eigenvalues, so the
Hoffman--Wielandt displacement term drops out.  The lemma
\texttt{gap\_of\_eigengap} supplies the Weyl step, converting a spectral gap of
$T$ (leading eigenvalues $\ge a$, trailing $\le b$) into the hybrid gap
$(a-b)-\eps$; YWS's $\Delta=\lambda_d-\lambda_{d+1}$ is the case $a-b$, and
\texttt{gap\_of\_rank\_floor} the case $a=\alpha$, $b=0$.  Their $d^{1/2}$
operator-norm branch is Corollary~\ref{cor:opnorm},
$\norm{\sinT}_{\F}\le d^{1/2}\eps/\gap$.

\section{What is not formalized}

The formalization covers the Frobenius $\sin\Theta$ theorem with a hybrid or
population gap.  The remaining DK-family results, and the machinery each would
require, are as follows.

\begin{itemize}
\item \emph{Operator-norm and general unitarily-invariant-norm $\sin\Theta$.} The
  dimension-free $\norm{\sinT}_{\op}\le\norm{S-T}_{\op}/\gap$, and its extension
  to every unitarily invariant norm (the full part~III statement).  The
  sum-of-squares route delivers the Frobenius norm; the operator-norm version
  needs the largest singular value of the residual, and the general one needs the
  theory of unitarily invariant norms (symmetric gauge functions, Ky Fan
  majorization), which Mathlib does not yet carry.
\item \emph{Davis's sharper total-rotation estimate.} The refinement that
  subtracts the eigenvalue displacement $\sum_i(\lambda_i-\lambda_i')^2$, and its
  $\norm{\mathcal C^{\perp}H}_{\F}$ form.  These rest on the canonical matching
  unitary, built from the polar factor of $P_j\widehat P_jP_j$, together with
  Davis's Schur--Horn eigenvalue-change lower bound.
\item \emph{The $\tan\Theta$, $\sin2\Theta$, and $\tan2\Theta$ theorems.} The
  other members of the DK family, each with its own gap condition and $2\times2$
  compression argument.
\item \emph{The YWS aligned-basis bound.} An orthogonal matrix $O$, built from the
  singular value decomposition of $\widehat V\tp V$, for which
  $\norm{\widehat V O-V}_{\F}\le 2^{3/2}\min\{\cdots\}/\Delta$.  The argument also
  uses the identity $2d-2\sum_k\cos\theta_k\le 2\sum_k\sin^2\theta_k$.
\item \emph{The YWS singular-vector extension.} The bound for singular subspaces
  of rectangular $A,\widehat A$, obtained by applying the symmetric result to
  $A\tp A$; it needs singular value perturbation and a Gram-perturbation bound on
  $\norm{\widehat A\tp\widehat A-A\tp A}$.
\item \emph{General-interval subspaces.} Spectral subspaces cut out by an
  arbitrary interval, with a two-sided gap, rather than the sorted leading block
  used here.
\end{itemize}

\section{Dictionary to the Lean source}

All declarations live in
\texttt{ForMathlib/Analysis/InnerProductSpace/DavisKahan.lean}
(namespace \texttt{ForMathlib}).

\begin{center}
\footnotesize
\setlength{\tabcolsep}{4pt}
\begin{tabular}{@{}p{0.42\linewidth}l@{}}
\toprule
Statement here & Lean name \\
\midrule
Cross-term identity \eqref{eq:cross}
  & \texttt{inner\_eigenvectorBasis\_map\_sub\_eigenvectorBasis} \\
Parseval identities (full sum / single row)
  & \texttt{sum\_sq\_norm\_inner\_\dots\_map\_sub\_eq}\,/\,\texttt{\dots\_eq\_row} \\
Dimension-factor step $\norm{S-T}_{\F}^2\le n\eps^2$
  & \texttt{sum\_norm\_eigenvectorBasis\_map\_sub\_sq\_le} \\
Ladder, cross-block rung $\norm{P(S-T)Q}_{\F}^2/\gap^2$
  & \texttt{sum\_cross\_\dots\_sq\_le\_offDiag} \\
Ladder, residual rung $\norm{(S-T)P}_{\F}^2/\gap^2$
  & \texttt{sum\_cross\_\dots\_sq\_le\_residual} \\
Ladder, sharp rung $\norm{S-T}_{\F}^2/\gap^2$ (Thm~\ref{thm:ladder})
  & \texttt{sum\_cross\_\dots\_sq\_le\_hilbertSchmidt} \\
Cor.~\ref{cor:opnorm}, $d^{1/2}$ branch $d\,\eps^2/\gap^2$
  & \texttt{sum\_cross\_\dots\_sq\_le\_opNorm} \\
Cor.~\ref{cor:opnorm}, crude $n\eps^2/\gap^2$
  & \texttt{sum\_cross\_\dots\_sq\_le} \\
Population-gap Weyl bridge $\Rightarrow\gap=(a-b)-\eps$
  & \texttt{gap\_of\_eigengap} \\
Projector identity $\norm{\widehat P-P}_{\F}^2=2\cdot\text{overlap}$
  & \texttt{sum\_norm\_sub\_starProjection\_span\_sq\_eq} \\
Projector form of Thm~\ref{thm:ladder}
  & \texttt{sum\_norm\_sub\_starProjection\_span\_sq\_le\_hilbertSchmidt} \\
Projector, crude $2n\eps^2/\gap^2$
  & \texttt{sum\_norm\_sub\_starProjection\_span\_sq\_le} \\
Weyl bridge, rank/floor $\Rightarrow\gap=\alpha/2$
  & \texttt{gap\_of\_rank\_floor} \\
Cor.~\ref{cor:floor}, rank/floor $4n\eps^2/\alpha^2$
  & \texttt{sum\_cross\_\dots\_sq\_le\_of\_rank\_floor} \\
\bottomrule
\end{tabular}
\end{center}

The projections $P,\widehat P$ are Mathlib's \texttt{Submodule.starProjection}
of the spans of the selected eigenvectors; the Frobenius norm on the left of the
projector bound is written in Lean as the basis sum
$\sum_k\norm{(\widehat P-P)u_k}^2=\norm{\widehat P-P}_{\F}^2$.

\begin{thebibliography}{9}
\bibitem{Davis1963}
Chandler Davis.
\newblock \href{https://doi.org/10.1016/0022-247X(63)90001-5}{The rotation of eigenvectors by a perturbation}.
\newblock \emph{J.\ Math.\ Anal.\ Appl.}\ \textbf{6} (1963), 159--173.

\bibitem{DavisKahan1970}
Chandler Davis and W.\ M.\ Kahan.
\newblock \href{https://doi.org/10.1137/0707001}{The rotation of eigenvectors by a perturbation.\ III}.
\newblock \emph{SIAM J.\ Numer.\ Anal.}\ \textbf{7} (1970), no.\ 1, 1--46.

\bibitem{YWS2015}
Yi Yu, Tengyao Wang, and Richard J.\ Samworth.
\newblock \href{https://doi.org/10.1093/biomet/asv008}{A useful variant of the Davis--Kahan theorem for statisticians}.
\newblock \emph{Biometrika} \textbf{102} (2015), no.\ 2, 315--323.
\newblock arXiv:\href{https://arxiv.org/abs/1405.0680}{1405.0680}.
\end{thebibliography}

\end{document}
